\section{The Finitary Realizer Normal Form compilation theorem}
\label{sec:compilation-theorem}

We now collect the preceding constructions into one theorem schema.  The statement is
metatheoretic: it takes fixed object syntax and local embedded derivations and returns a
new embedded derivation.  It does not assert that \(\iPRA\) internally quantifies over
all such graphs.

\subsection{FRNF graph}

\begin{definition}[Certified FRNF graph]
A certified \(\FRNF\) graph consists of the following finite data.

\begin{enumerate}[label=\textnormal{G\arabic*.},leftmargin=3.4em]
\item A finite acyclic directed graph with a designated final node and fixed
 topological order.
\item For each input, a finite-observation presentation: exact data, regular names,
 finite clouds, finite nets, or certified setoids.
\item For every edge, primitive-recursive computation and justification horizon
 transformers with an embedded proof that read demand is bounded by justification
 demand.
\item For every node, fixed primitive-recursive syntax and an object contract formula.
\item Each computational node is one of: exact transducer, bounded search, clocked
 loop, finite relational viability pass, stable observer, gap collapse, boundary node,
 or finalizer.  Compound domain routines are expanded into these forms.
\item Every clocked loop has a bounded/QF invariant and embedded base, step, and finish
 derivations.
\item Every bounded formula occurring in a local invariant has a characteristic term
 and the embedded Booleanization derivations of \cref{thm:booleanization}.
\item Every node has an embedded local derivation of its contract from its declared
 predecessor contracts and permanent boundary formulas.
\item Every relational output consumed deterministically is either internal to a
 finite viability computation or followed by an embedded stability/gap derivation.
\item Every non-finite mathematical input---modulus, height bound, separation, finite
 net, zero-free estimate, covering list---is either produced by a preceding certified
 node or appears explicitly among the boundary formulas.
\end{enumerate}
\end{definition}

The definition is intentionally syntactic.  Host-semantic theorems may accompany the
graph, but they are not among the premises needed to construct the object proof.

\subsection{Compilation theorem}

\begin{theorem}[FRNF compilation]
\label{thm:frnf-compilation}
Let \(G\) be a certified \(\FRNF\) graph and let \(d\) be a final demand term.  There is
a metatheoretically effective construction that produces:

\begin{enumerate}[label=(\roman*),leftmargin=2.4em]
\item primitive-recursive terms coding the least dual-horizon demand closure of
\((G,d)\);
\item finite input-table and trace terms sufficient for every computation and proof in
\(G\);
\item an embedded derivation
\[
   \Gamma_G\vdash_{\PRAAx}\operatorname{Goal}_G(d),
\]
where \(\Gamma_G\) is exactly the list of undischarged boundary formulas; and
\item by the existing transfer theorem, a derivation of the translated goal in the
native scalar core.
\end{enumerate}
No induction beyond quantifier-free induction is introduced by the construction.
\end{theorem}

\begin{proof}
We describe the compiler and verify each stage.

\paragraph{1. Static specialization.}
Replace every externally fixed routine by concrete primitive-recursive syntax and its
defining equations.  If a routine was supplied as a finite table or as a code in a
fixed restricted language, instantiate the corresponding table lookup or bounded
interpreter.  No total universal evaluator is introduced.

\paragraph{2. Reverse dual-demand pass.}
Apply \cref{thm:demand-closure} to the finite graph at final demand \(d\).  This yields
for each node/input the least read and justification bounds.  The pass itself is a
bounded primitive recursion on the encoded reverse topological list; its correctness
has a QF invariant.

\paragraph{3. Finite materialization.}
Construct codes for every input row, validity instance, approximation, candidate, and
edge relation at indices within the justification closure.  Exact inputs are constant
tables.  For regular names, the global name hypotheses are instantiated only at the
finitely many requested indices.  For relational inputs, materialize the bounded
candidate tables.

\paragraph{4. Booleanization.}
Apply \cref{thm:booleanization} to every bounded validity, range, table, and edge
condition.  Local state predicates become equations \(\chi=1\), and finite universal
conditions become bounded-product equations.  Thus all loop invariants are
quantifier-free.

\paragraph{5. Local computational nodes.}
Exact transducers use defining equations.  Bounded searches use
\cref{thm:bounded-search}.  Majorized descents use
\cref{thm:majorized-descent}.  Finite relational nodes use
\cref{thm:backward-viability}.  Each adaptive routine has been defunctionalized into a
clocked transition system, so \cref{thm:trace-internalization} constructs its embedded
contract from the supplied base, step, and terminal derivations.

\paragraph{6. Choice erasure.}
When a downstream observable consumes relational candidates, use the supplied embedded
stability theorem and \cref{thm:stable-observer}.  If an exact discrete label is
requested, use the explicit gap data and \cref{thm:gap-collapse}.  No arbitrary exact
branch choice is generated.

\paragraph{7. Topological proof composition.}
Process graph nodes in topological order.  By assumption, each local contract is
derivable from declared predecessors and \(\Gamma_G\).  Apply
\cref{thm:proof-schedule} at every node.  Induction on the finite list yields a
derivation of the final goal from \(\Gamma_G\).

\paragraph{8. Native export.}
Apply \code{pra\_iff\_native} to the completed ordinary-iPRA derivation.  This stage
changes presentation but not mathematical assumptions.

Every induction introduced above is on a numerical clock/table index with a
Booleanized state equation.  Hence it is quantifier-free induction.  All other logical
steps are ordinary intuitionistic deduction and equality reasoning.  The only
undischarged hypotheses are the boundary formulas because all generated component
contracts have been cut away.
\end{proof}

\subsection{Relative expressiveness}

The theorem is not a completeness theorem for all constructive mathematics.  It is a
normal-form theorem for proofs already equipped with primitive-recursive finite
realizers and quantitative interfaces.  Within that class, the node vocabulary is
broad: every fixed primitive recursion is a clocked loop; every bounded existential is
least search; every finite nondeterministic process is a layered relation; every finite
DAG composes by the schedule.  Thus the remaining expressiveness questions are mainly
about whether the mathematics supplies finite bounds and stable representations.

\subsection{Conservativity of the compiler}

The compiler adds no axioms to \(\iPRA\).  Its generic Lean definitions live in the
metatheory and build terms of the inductive relation \(\Prf\,\PRAAx\).  Trust rests on
the Lean kernel, the definition of \(\Prf\), and the audited local constructors.  A
future executable compiler may be large or use sophisticated metaprogramming without
changing the object-theory theorem, provided the generated proof term checks.

\subsection{The main engineering gap}

The logical schedule and QF loop combinators are short.  The substantial remaining
implementation is a reusable object arithmetic layer for:

\begin{itemize}
\item generic bounded-formula reification and substitution;
\item encoded arrays and finite folds;
\item graph demand tables;
\item viability masks and path extraction; and
\item rational error budgets and metric inequalities.
\end{itemize}

The checked \code{IPRADerivations} library already contains most low-level ingredients.
The task is integration and abstraction, not a new foundational theory.
